\chapter{Tools}\label{ann:outils}

Every \mes{} or \calc{} number of the thesis comes out of a committed tool. This appendix lists them, sorted by study, each with the first sentence of its own description, quoted as is (the tools are written in French, as is the whole repository). The table is generated by \fichier{memoire/outils/annexes.py} from \barre{\texttt{git ls-files}} \texttt{git ls-tree}: it is not written by hand and \barre{follows the repository} describes the repository at the publication of this volume, at commit \texttt{b8af154}. Corrected on 2026-10-29: the tool read the repository of the day, which kept growing after this volume; it now reads this commit, and the table stays that of the published volume.

\section{The thesis toolchain}

\begin{itemize}
\item \fichier{memoire/construire.sh} (and \fichier{memoire/Makefile} for those who have \texttt{make}): regenerates the figure data and the appendix tables, then builds both PDFs with \texttt{latexmk} (lualatex, biber); it fails at the first error, at any missing reference or citation, and at any entry biber cannot find.
\item \fichier{memoire/outils/donnees_figures.py}: extracts from each committed log the CSV of one figure, with no number written by hand (the only constants are the datasheet ceilings, each with its derivation).
\item \fichier{memoire/outils/annexes.py}: the tables of this appendix and the next.
\item \fichier{memoire/outils/bib_etudes.py}: the bibliographies of the studies, in biblatex.
\end{itemize}

\section{The tools of the studies}

\input{\Racine/annexes/outils}
